\section{An exact one-edge escape from the coherent family}
\label{sec:one-edge}

The pinned source contains a certified counterexample to coherent-only
normal-disc discovery. Start with the two-tetrahedron layered solid torus,
then perform six legal $2$--$3$ moves, choosing the following
$(\text{tetrahedron},\text{face})$ in the current labelling:
\[
                 (1,2),(0,1),(2,2),(3,1),(3,2),(6,0).
\]
The full source, intermediate face pairings, and move certificates are in
\path{synthesis/data/coherent-obstruction-certificate.json} of the
pinned repository. Its final triangulation $T_*$ has eight tetrahedra,
one vertex and ten edges. Its Regina isomorphism signature is
\begin{center}\texttt{iHvrQAkbgfeffhhasssvcaj}.\end{center}
Legal Pachner moves preserve the
underlying solid torus; the independently replayed source history is
part of the evidence, not an inference from a successful disc query.

\subsection{The coherent surface and the new vertex disc}

The following table gives the primitive coherent surface $C$ and the new
meridian $D$. Coordinates in a row are
$(t_0,t_1,t_2,t_3;q_{01\mid23},q_{02\mid13},q_{03\mid12})$.
The semicolon separates triangle and quadrilateral coordinates.
\begin{table}[htbp]
\centering
\begin{tabular}{@{}c cc@{}}
\toprule
Tetrahedron & Coherent $C$ & Vertex disc $D$\\
\midrule
0 & $(2,2,0,0;0,0,2)$ & $(2,2,0,0;0,0,2)$\\
1 & $(0,0,2,2;0,2,0)$ & $(0,0,2,2;0,2,0)$\\
2 & $(0,0,1,2;0,0,2)$ & $(0,0,1,2;0,0,2)$\\
3 & $(0,2,0,2;0,0,1)$ & $(0,2,0,2;0,0,1)$\\
4 & $(2,0,0,2;0,1,0)$ & $(2,0,0,2;0,1,0)$\\
5 & $(0,1,0,1;3,0,0)$ & $(2,3,0,1;1,0,0)$\\
6 & $(0,1,2,0;2,0,0)$ & $(2,3,2,0;0,0,0)$\\
7 & $(0,1,0,1;2,0,0)$ & $(2,3,0,1;0,0,0)$\\
\bottomrule
\end{tabular}
\caption{An exact 47-piece meridian outside the primitive coherent family.
All coordinates are small integers; the accompanying certificates bind
them to the full face-pairing source.}
\label{tab:coordinates}
\end{table}

In increasing validator edge-root order, the weights are
\begin{center}
\small
\begin{tabular}{@{}c rrrrrrrrrr@{}}
\toprule
Edge root &0&3&4&6&13&15&18&22&24&30\\
\midrule
$w(C)=b$ &6&2&4&2&3&1&3&5&3&1\\
$w(D)$   &6&2&4&2&3&1&3&5&3&5\\
\bottomrule
\end{tabular}
\end{center}
The boundary edge roots are $3,13,22$, so the changed edge $30$ is
interior. Its cancellation multiplicity is $z_{30}=(5-1)/2=2$.
All other cancellation multiplicities vanish.

The exact certificates give $C$ 41 pieces, Euler characteristic $-3$,
one connected orientable component and one boundary curve. It has genus
two. The new $D$ has 47 pieces, Euler characteristic one and one essential
disc component. Its standard support contains 25 positive coordinates;
the supported matching matrix has rank 24. Thus the ray is a standard
vertex ray, not only a ray made extreme by the selected weight space.
The primitive cocycle lattice has rank one, so every primitive coherent
surface comes from $c$ or $-c$ and has exactly the same unsigned vector.
A zero-cancellation meridian would therefore equal $C$, which is
impossible. Consequently
\[
                             \kappaV(T_*)=1.
\]

The previously retained 51-piece disc changes two edge weights, at roots
$15$ and $18$, each by four. The new disc improves that explicit witness
and uses fewer cancellation edges. It does \emph{not} enlarge coverage
beyond the existing complete search in the coherent surface's full
quadrilateral sector: $D$ belongs to that sector, with two of its allowed
quadrilateral coordinates becoming zero. The new contribution is a
complete parameterization and a less expensive way to discover the disc.

\subsection{The whole two-dimensional cone is explicit}

Consider the affine slice of $W_{\{30\}}$ in which the coherent coefficient
is one:
\[
                         w(s)=b+s\mathbf e_{30}.
\]
Only tetrahedra 5, 6 and 7 change under the local reconstruction. In each,
the coordinates $t_0,t_1$ increase by $s/2$; their occupied $q_0$
coordinates become respectively
\[
                         3-s/2,\quad 2-s/2,\quad 2-s/2.
\]
All remaining coordinates are fixed. These vectors are admissible exactly
for $0\le s\le4$. Negative $s$ would make $t_0$ negative, while $s>4$
fails the opposite-sum maximum condition in tetrahedra 6 and 7.
Nonnegativity on edge $0$ forces the coherent coefficient $\lambda\ge0$.
There are no admissible nonzero rays with coherent coefficient zero:
putting weight only on edge $30$ gives a negative face-corner expression.
Thus the admissible cone in this weight space has just the two boundary
rays $C$ and $D$.

For $0\le s\le4$, direct cell counting gives
\[
                   \chi(w(s))=-3+s,\qquad
                   \#\text{pieces}=41+\frac{3s}{2}.
\]
At $s=0,2,4$ the vectors are integral. Independent full topology replay
finds a connected orientable surface with one boundary curve in every
case, of genus $2,1,0$, respectively. The geometry behind the search is
therefore visible: the endpoint $s=4$ removes the remaining handle
obstruction while changing only one interior edge weight. The genus-one
intermediate vector is not an extreme ray of this two-dimensional cone.

The complete, untruncated arrangement enumeration on this source gives:
\begin{center}
\begin{tabular}{@{}r r r r r@{}}
\toprule
Cancellation radius & Weight spaces & Intersections & Admissible rays & Disc vectors\\
\midrule
0 & 1 & 1 & 1 & 0\\
1 & 11 & 73 & 2 & 1\\
2 & 56 & 5,634 & 5 & 2\\
\bottomrule
\end{tabular}
\end{center}
All five rays found at radius two happen to be standard vertex rays.
This observation is specific to the source; the general algorithm may
emit rays that are extreme only in a restricted space. The counts here
include complete enumeration and therefore differ from early-return
search counters in the benchmark.

\section{An infinite dense family with one cancellation edge}
\label{sec:layering}

The following elementary extension lemma turns the certified eight-tetrahedron
example into an infinite family.  It also explains why large coordinates
and quadrilaterals in almost every tetrahedron need not enlarge the new
discovery parameter.

\subsection{A support-rigid coherent extension}

Let a tetrahedron $\Delta=0123$ be attached along the faces $012$ and
$013$, so that its old edges are all the edges except $23$.  For real
heights $h_0,h_1,h_2,h_3$, the \emph{coherent local vector} is obtained
by taking the level sections of the affine height function.  If
$h_{v_0}\leq h_{v_1}\leq h_{v_2}\leq h_{v_3}$, its potentially nonzero
coordinates are the triangle at $v_0$, the quadrilateral separating
$\{v_0,v_1\}$ from $\{v_2,v_3\}$, and the triangle at $v_3$, with
coefficients respectively
\[
 h_{v_1}-h_{v_0},\qquad h_{v_2}-h_{v_1},\qquad
 h_{v_3}-h_{v_2}.
\]
Tied heights cause zero coordinates; the statement includes those cases.

\begin{lemma}[Injectivity on the coherent support]
The five old edge weights determine a normal vector supported in the
positive support of this coherent local vector uniquely.
\end{lemma}

\begin{proof}
The supported coordinates are successive height gaps in the displayed
fixed ordering.  Their edge weights are signed height differences, with
the signs fixed by that ordering.  Knowing the weights on the connected
graph $K_4\setminus\{23\}$ determines the four heights up to a common
additive constant, and hence determines every height gap.  This argument
is linear and remains valid when some of the gaps are constrained to zero.
\end{proof}

\begin{proposition}[Preservation of a standard vertex ray]
Suppose a standard vertex normal vector $x$ on the old triangulation
extends across $\Delta$ by adding a coherent local vector $u$.  Then
$x'=(x,u)$ is again a standard vertex vector after primitive normalization.
If $x$ is primitive integral and $u$ integral, no normalization is needed.
\end{proposition}

\begin{proof}
In any nonnegative matching decomposition $x'=y+z$, restriction to the old
tetrahedra gives a decomposition of $x$.  Its vertex property gives
$y_{\mathrm{old}}=a x$ and $z_{\mathrm{old}}=(1-a)x$ for some
$0\leq a\leq1$.  Matching on the attaching faces fixes the five old edge
weights of $y_\Delta$ to be $a$ times those of $u$.  Nonnegativity forces
$y_\Delta$ to be supported in the positive support of $u$.  The preceding
lemma gives $y_\Delta=au$.  Thus $y=ax'$ and the ray is extreme.  An
integral extension of an integral vector with coordinate gcd one still
has coordinate gcd one.
\end{proof}

\subsection{Euler characteristic under a torus diagonal flip}

Assume the attaching faces are the two triangles of a one-vertex torus,
and that the local heights obey the parallelogram identity
\[
 h_0+h_1=h_2+h_3.                                      \tag{P}
\]
This identity is the cocycle compatibility condition on the four sides
of the lifted quadrilateral.  Adding $\Delta$ performs the boundary
diagonal flip and preserves the underlying solid torus.

For a finite set $A$ of real heights, write
$\operatorname{span}(A)=\max A-\min A$.  The change in the normal
surface Euler characteristic is
\begin{align*}
 \Delta\chi={}& |h_2-h_3|
 -\operatorname{span}(h_0,h_2,h_3)
 -\operatorname{span}(h_1,h_2,h_3)\\
 &+\operatorname{span}(h_0,h_1,h_2,h_3).                \tag{E}
\end{align*}
Indeed, the new edge $23$ contributes the first term to the normal
surface's zero-cells.  The two new boundary faces contribute the two
subtracted arc counts.  The coherent pieces in the new tetrahedron
contribute the final term.  All old cells retain their previous counts.

\begin{lemma}[Zero Euler correction]
Under (P), the right hand side of (E) is zero.
\end{lemma}

\begin{proof}
Put $a=\min(h_2,h_3)$ and $b=\max(h_2,h_3)$.  Since
$h_0+h_1=a+b$, either both $h_0,h_1$ lie in $[a,b]$, or one lies
$\delta\geq0$ below $a$ and the other lies $\delta$ above $b$.
In the first case, all three spans in (E) equal $b-a$.
In the second case, the two subtracted spans equal $b-a+\delta$
and the last equals $b-a+2\delta$.  Both cases give zero.
\end{proof}

If $x$ is a primitive vertex meridian disc whose weights on the three
boundary edges are coherent for the primitive cocycle, this extension
is again a meridian disc.  The extension is a primitive vertex surface,
so it is connected: otherwise its distinct component vectors would
decompose its ray, and primitivity excludes multiple copies.  Its Euler
characteristic is one by the lemma.  Its new boundary is the coherent
curve dual to the same primitive boundary cocycle, which is a single
essential curve on the two-triangle torus.  Therefore the connected
surface is an essential disc.  Equivalently, the boundary statement
follows by writing the two-triangle torus as an affine parallelogram
quotient: a primitive integral linear functional has a connected
essential level curve.

The coherent surface for the primitive cocycle also preserves its genus
and number of components in this construction.  Its coorientation
extends across the tetrahedron.  Every new coherent normal piece meets
at least one attaching face: their two height intervals overlap on the
height interval of edge $01$ and their union is the complete height
interval of $\Delta$.  Consequently a connected old coherent surface
remains connected.  The boundary remains one primitive coherent curve,
and the Euler correction is zero.

\subsection{The explicit Fibonacci family}

Let $T_0=T_*$ be the certified eight-tetrahedron solid torus from
Section~\ref{sec:one-edge}. Its primitive coherent surface
$C_0$ has $41$ normal pieces, Euler characteristic $-3$, and one boundary
component, so it has genus two.  Let $D_0$ be the new $47$-piece standard
vertex meridian given in the accompanying exact certificate.  It differs
from the absolute cocycle weights only at the interior edge numbered
$30$, where $w_{D_0}(30)=|c(30)|+4=5$.  Thus its cancellation support
has exactly one edge.  The boundary weights are $2,3,5$.

At each step layer a tetrahedron along the boundary edge of smallest
weight.  With $F_1=F_2=1$, the $j$th added tetrahedron, $j\geq1$, has
coherent heights, after a local relabelling and perhaps reversing all
signs,
\[
 (h_0,h_1,h_2,h_3)
 =\bigl(0,F_{j+2},-F_{j+3},F_{j+4}\bigr).
\]
The identity $F_{j+4}=F_{j+3}+F_{j+2}$ gives (P).  The old diagonal
has weight $F_{j+2}$ and the new diagonal has weight $F_{j+5}$.
The two other boundary edge weights are unchanged.  Hence the three
boundary weights after $r$ layers are
\[
 F_{r+3},\quad F_{r+4},\quad F_{r+5}.
\]
The new local triangle coefficients are $F_{j+3}$ at the height
minimum and maximum, and its quadrilateral coefficient is $F_{j+2}$.
It contributes $2F_{j+3}+F_{j+2}=F_{j+5}$ normal pieces.

\begin{theorem}[Dense one-cancellation family]
\label{thm:dense-family}
For every $r\geq0$, the preceding construction gives a one-vertex
solid-torus triangulation $T_r$ with $t=r+8$ tetrahedra and a primitive
standard vertex meridian $D_r$ such that:
\begin{enumerate}
\item its cancellation support is exactly the original interior edge $30$;
\item its quadrilateral support has size $r+6=t-2$;
\item its number of normal pieces is $F_{r+7}+34$;
\item the primitive coherent surface $C_r$ is connected, has one boundary
component, genus two, and $F_{r+7}+28$ pieces.
\end{enumerate}
In particular, dense quadrilateral support, exponentially large normal
coordinates, and persistent failure of the coherent seed are compatible
with a cancellation-support discovery parameter equal to one.
\end{theorem}

\begin{proof}
The base assertions are established by the supplied exact certificates
and Table~\ref{tab:coordinates}. All four vertices of each new tetrahedron
belong to attaching faces, so the source retains one vertex. The old
cocycle coordinates have gcd one and are unchanged, so the extended
cocycle stays primitive.
The support-injectivity proposition preserves the standard vertex ray,
and the Euler and boundary arguments preserve the meridian disc and the
genus-two coherent surface.  All new coordinates are coherent for the
extended cocycle; the old defect edge is interior and never changes.
Thus the cancellation support stays exactly $\{30\}$. Since $C_r$ remains a genus-two surface, no meridian has empty cancellation support; hence $\kappaV(T_r)=1$. Each new
tetrahedron has a positive quadrilateral coefficient, so the support
increases by one at each step.  Finally,
\[
 \sum_{j=1}^r F_{j+5}=F_{r+7}-F_7=F_{r+7}-13.
\]
Adding $47$ or $41$ yields the stated piece counts.  The largest
coordinate has $\Theta(r)$ bits even though its value grows exponentially.
\end{proof}

This theorem concerns the discovery parameter.  The base disc lies in
the same full quadrilateral sector as the coherent surface, so the
family does not prove a coverage improvement over an existing complete
search supplied with that sector.  The new theorem keeps a low-dimensional
edge-weight restriction through ray discovery, independently of the
number of tetrahedra with quadrilaterals.


Exponential coordinate growth in layered solid tori is established in
normal-surface complexity theory; Burton gives an explicit Fibonacci
vertex-disc construction \cite{burton-complexity}. The distinguishing
feature here is its persistence together with a genus-two coherent
obstruction and exactly one cancellation edge. The table below gives
some exact sizes; the saved benchmark also reaches $r=64$.
\begin{center}
\begin{tabular}{@{}r r r r@{}}
\toprule
Layers $r$ & Tetrahedra & Disc pieces & Occupied quadrilateral types\\
\midrule
0 & 8 & 47 & 6\\
1 & 9 & 55 & 7\\
2 & 10 & 68 & 8\\
4 & 12 & 123 & 10\\
8 & 16 & 644 & 14\\
16 & 24 & 28,691 & 22\\
32 & 40 & 63,246,020 & 38\\
64 & 72 & 308,061,521,170,163 & 70\\
\bottomrule
\end{tabular}
\end{center}

\section{A linear cancellation obstruction for arbitrary sources}
\label{sec:whitehead}

A logarithmic cancellation bound cannot hold for all finite
triangulations. The obstruction below uses multiple vertices; it is
therefore separate from the one-vertex theorem and does not resolve its
most interesting remaining geometric question.

\begin{theorem}[Many forced cancellation edges]
\label{thm:whitehead}
There are generalized triangulations $U_r$ of a solid torus with
$t_r=6r+O(1)$ tetrahedra such that every oriented meridian disc transverse
to their one-skeleton has at least $r$ cancellation edges. In particular
$\kappaV(U_r)\ge r=\Omega(t_r)$.
\end{theorem}
\begin{proof}
Choose a fixed Whitehead curve $J$ in a solid torus $V$. Its winding
number is zero, while its geometric index is two: every meridian disc
of $V$ meets $J$ at least twice. These properties are established, for
example, by the geometric-index calculations in Section 5 and Figure 3(a)
of Andrist--Garity--Repov\v{s}--Wright \cite{geometric-index}.

Choose a fixed tubular neighbourhood $N(J)\cong S^1\times D^2$.
Triangulate $D^2$ with three fixed boundary vertices and $r$ interior
vertices $p_1,\ldots,p_r$. Repeated insertion of an interior vertex into
one triangle gives exactly $2r+1$ triangles and leaves the boundary
triangulation unchanged. Put the one-vertex, one-edge $\Delta$-complex
structure on $S^1$. Give the disc vertices a global order, retaining the
boundary order, and triangulate every product triangular prism by the
compatible three-tetrahedron staircase subdivision. The resulting
product has $3(2r+1)$ tetrahedra. Each
\[
                            e_j=S^1\times\{p_j\}
\]
is one global loop edge, and the $r$ loop edges are distinct and disjoint.
The triangulation induced on $\partial N(J)$ is independent of $r$.

Glue to a fixed triangulation of
$V\setminus\operatorname{int}N(J)$. If their boundary triangulations
need reconciliation, use a fixed triangulated collar. It does not touch
the interior loop edges. The total tetrahedron count is $6r+O(1)$ and
the underlying manifold is the solid torus $V$.

Every $e_j$ is parallel to $J$ in its fixed neighbourhood, hence has
winding number zero and geometric index two in $V$. For every oriented
meridian $D$,
\[
                         c_D(e_j)=0,\qquad w_D(e_j)\ge2.
\]
Equation~\eqref{eq:defect} makes each $e_j$ a cancellation edge. This
argument does not depend on a choice of vertex potential: every
coboundary evaluates to zero on a loop edge. There are $r$ such edges,
which proves the bound for all meridians and hence for vertex meridians.
\end{proof}

The theorem is a representation obstruction, not a lower bound for
unknot recognition. Its source has $\Theta(r)$ vertices, and the theorem
does not prove that these obstructions survive a controlled reduction to
one vertex. The product construction deliberately makes the loops
separate vertices so that their triangulation cost and cancellation
contributions are transparent.

A successful global approach must therefore choose its sources. It may
need certified retriangulations that remove such repeated gadgets, a
parameter measuring independent obstructions rather than literal edge
support, or a family of alternative exterior presentations. A bound on
one convenient triangulation chosen by an unspecified topological
existence argument would still be insufficient: both discovery of that
source and the cost of certifying its relation to the diagram belong in
the complexity theorem.
